\documentclass[12pt]{article}
\usepackage[utf8]{inputenc}
\usepackage[T1]{fontenc}
\usepackage{hyperref}
\title{Fang Language Resources for K-12 Curriculum Development}
\author{Compiled by Claude Sonnet 5.5 (claude-sonnet-5-5) for Andr\'as Kornai}
\date{}
\begin{document}
\maketitle

\section{Introduction}
Fang (ISO 639-3 \texttt{fan}, Glottocode fang1246) is a Bantu language spoken in Equatorial Guinea, northern Gabon and southern Cameroon. Its dialects include Ntumu, Okak, Mek\^e, Mva\"i, Atsi and Nzaman. Several orthographies are in use. The searches found reference and scripture-audio material but no school textbooks or curriculum materials in Fang.

\section{Primary Educational Resources}
None found. No Fang primers or school textbooks appeared in search results.

\section{Linguistic and Reference Materials}
\begin{itemize}
\item Wikipedia, ``Fang language'': overview and dialects. \url{https://en.wikipedia.org/wiki/Fang_language}
\item Omniglot page on Fang writing; notes several orthographies in use. \url{https://omniglot.com/writing/fang.htm}
\item Royal Museum for Central Africa project compiling a bilingual Fang-French / French-Fang dictionary (about 15,000 entries planned, based on Galley 1964 and the Atsi dialect). \url{https://www.africamuseum.be/en/research/discover/projects/prj_detail?prjid=214}
\item Glottolog reference record for Amado Osorio Zabala (1887), a Fang-Spanish vocabulary. \url{https://glottolog.org/resource/reference/id/66800}
\end{itemize}

\section{Religious and Community Resources}
\begin{itemize}
\item Joshua Project profile of the Fang in Gabon; its Scripture-status statements conflict across pages, so they should be checked against another source. \url{https://www.joshuaproject.net/index.php/people_groups/11187/GB}
\end{itemize}

\section{Digital and Audio Resources}
\begin{itemize}
\item Global Recordings Network page for Fang. \url{https://globalrecordings.net/en/language/fan}
\item Dramatized audio New Testament in Fang (Megavoice listing). \url{https://megavoice.com/media-cloud/m0b6415-new-testament-dramatized-fan-fang-audio-bible}
\end{itemize}

\section{Curriculum Development Context}
Fang is widely spoken but, as far as these searches show, lacks published school materials. The most useful starting points are the dictionary project, the older vocabularies and the audio scripture, which would need verification and a choice among orthographies and dialects before curriculum use.

\end{document}